\providecommand{\CRevisionRound}{2}
\documentclass[9pt]{extarticle}
\usepackage[a4paper,margin=0.68in]{geometry}
\usepackage{fontspec}
\setmainfont{Latin Modern Roman}
\setsansfont{Latin Modern Sans}
\setmonofont{Latin Modern Mono}
\usepackage{amsmath,amssymb,booktabs,microtype,enumitem,xcolor,hyperref}
\hypersetup{colorlinks=true,linkcolor=blue!45!black,citecolor=blue!45!black,urlcolor=blue!45!black}
\setlist{nosep,leftmargin=1.25em}
\newcommand{\TV}{\mathrm{TV}}
\newcommand{\A}{\mathsf A}
\title{The Base Semigroup of Holte's Carries Chain:\\An Exact Route-A Certificate}
\author{Route-A structural certificate HCS-C194}
\date{}

\begin{document}
\maketitle
\begin{abstract}
Adding $n$ independent base-$b$ digit columns produces a Markov chain on the
carry states $\{0,\ldots,n-1\}$.  We source-lock Holte's all-parameter
transition formula, mixed-radix semigroup $P_aP_b=P_{ab}$, common eigenvectors,
simple spectrum $1,b^{-1},\ldots,b^{-(n-1)}$, and Eulerian stationary law.
The diagonalization closes every power trace, finite determinant, and an exact
common-projector convergence expansion; a sourced total-variation estimate is
kept separate from that algebraic deduction.  An exact rational census checks
72 matrices through independent algorithms but is only regression evidence.
An ownership and control ledger prevents classical results or prime-tagged
fixtures from being recast as package novelty.  Ordinary
positional addition earns a weak arithmetic relation; prime and composite
bases obey the same theorem, so no rational-prime orbit semantics, target
divisor, or natural target quantization follows.
\end{abstract}
\noindent\textbf{Keywords:} carries; positional addition; Markov chain;
Eulerian numbers; matrix semigroup; Route A.

\section{The frozen digit-column dynamics}
Fix $n\ge1$ and $b\ge2$.  From a carry-in $i\in X_n=\{0,\ldots,n-1\}$,
sample $d_1,\ldots,d_n$ independently and uniformly from
$\{0,\ldots,b-1\}$.  One clock step returns the unique $j$ such that
\[
 jb\le i+d_1+\cdots+d_n<(j+1)b.
\]
Introducing the output remainder as a slack digit gives Holte's coefficient
formula~\cite[Theorem~1]{Holte97}
\begin{equation}\label{eq:transition}
 P_b(i,j)=b^{-n}[x^{(j+1)b-1-i}]
 (1+x+\cdots+x^{b-1})^{n+1}.
\end{equation}
The chain is irreducible and aperiodic.  The case $n=1$ is the one-state
identity, retained as the exact boundary.

\section{Base semigroup and complete spectrum}
Write a base-$ab$ digit uniquely as $x+ay$, with $0\le x<a$ and
$0\le y<b$.  Applying the carry first to the base-$a$ column and then to the
base-$b$ column gives the same output as one base-$ab$ column.  Therefore
\begin{equation}\label{eq:semigroup}
 P_aP_b=P_{ab},\qquad P_b^r=P_{b^r}.
\end{equation}
Holte's Theorem~3 gives a matrix $V_n$, independent of $b$, such that
\begin{equation}\label{eq:diagonal}
 V_nP_bV_n^{-1}=\operatorname{diag}
 (1,b^{-1},\ldots,b^{-(n-1)}).
\end{equation}
The eigenvalues are distinct.  The first row of $V_n$ consists of Eulerian
numbers, hence the stationary probability is
\[
 \pi_n(j)=\frac{\A(n,j)}{n!},\qquad0\le j<n.
\]
It follows immediately that
\begin{align}
 \operatorname{tr}(P_b^r)&=\sum_{k=0}^{n-1}b^{-rk},\label{eq:trace}\\
 \det(I-zP_b)&=\prod_{k=0}^{n-1}(1-zb^{-k}),\label{eq:det}\\
 \chi_{P_b}(t)&=\prod_{k=0}^{n-1}(t-b^{-k}).\label{eq:char}
\end{align}
These are finite Markov-operator identities, not an Artin--Mazur or target
determinant.

\section{Common projectors and convergence}
Let $E_0,\ldots,E_{n-1}$ be the spectral projectors determined by the
base-independent matrix $V_n$.  They are independent of $b$, and
\begin{equation}\label{eq:projectors}
 P_b^r-E_0=\sum_{k=1}^{n-1}b^{-rk}E_k.
\end{equation}
Here $E_0$ maps every row distribution to $\pi_n$.  Equation
\eqref{eq:projectors} is an exact identity, not merely an asymptotic statement;
in every fixed matrix norm it gives geometric convergence governed by
$b^{-r}$ unless the $E_1$ component vanishes.

Diaconis and Fulman identify the carries law with a descent marginal of
repeated riffle shuffles and analyze convergence~\cite{DF09}.  Their
Theorem~3.3 supplies, for $n\ge3$, every start $i$, and $r\ge0$,
\begin{equation}\label{eq:tv}
 \lVert P_b^r(i,\cdot)-\pi_n\rVert_{\TV}
 \le \frac{(n-1)/2+i}{b^r}.
\end{equation}
We attribute this bound rather than presenting it as a package theorem.
For $n=1$ the chain is already stationary; $n=2$ follows directly from the
two eigenvalues $1,b^{-1}$ and is not folded into the stated $n\ge3$ source
locator.

\section{Exact regression certificate}
The producer reconstructs digit-sum coefficients by convolution.  A separate
checker uses slack-variable inclusion--exclusion, enumerates Eulerian numbers
from permutations, and obtains characteristic coefficients through the
Faddeev--LeVerrier algorithm.  SymPy independently checks matrix polynomials,
eigenspace dimensions, stationary equations and traces.

The frozen census contains 72 cases ($1\le n\le8$, $2\le b\le10$), 1,836
transition cells, 392 base-semigroup tuples and 96 power-identity tuples.
Bases $2,3,5,7$ contribute 32 prime-tagged cases; $4,6,8,9,10$ contribute 40
composite-tagged controls.  Enumeration tests code only: the all-parameter
quantifiers belong to Holte's theorem.

The independent checker passes 24,602 assertions, while the separate SymPy
oracle passes 14,248 exact checks.  Isolated replay reproduces all 537,471
evidence bytes.  A hostile suite rejects 159 repaired-hash semantic attacks and
one stale-hash attack.  These figures are released outputs, not estimates.

\section{Ownership and control ledger}
\begin{table}[ht]
\centering\small
\begin{tabular}{@{}p{0.25\linewidth}p{0.31\linewidth}p{0.34\linewidth}@{}}
\toprule
Statement & Owner or derivation & Release ceiling \\
\midrule
transition coefficient & Holte Theorem 1 & no proof by finite census \\
common eigenbasis and spectrum & Holte Theorem 3 & no novelty claim \\
right eigenvectors & Holte Theorem 4 & not a package diagonalization \\
descent marginal and TV analysis & Diaconis--Fulman Theorems 1.1, 3.1, 3.3 & no shuffle theorem is split off \\
trace, determinant, projectors & elementary finite-dimensional deduction & no target divisor \\
prime/composite cases & exact regression controls & no prime privilege \\
\bottomrule
\end{tabular}
\caption{Classical ownership, package deductions, and claim ceilings.}
\end{table}

The stable JSTOR identifier for Holte's paper is
\href{https://doi.org/10.2307/2974981}{doi:10.2307/2974981}; the publisher DOI
used in the registry is printed below.  The two-source population, theorem
locators, scope literal, and attribution status are exact-matched by the
checker.  The package claims neither global literature priority nor external
peer review.

\section{Route-A verdict}
Positional addition is genuinely arithmetic, but its base-power semigroup is
not a rational-prime primitive-orbit repetition law and supplies neither a
$\log p$ clock nor arithmetic weights.  The frozen object is a stochastic matrix rather than a deterministic
primitive-orbit map.  Its finite determinant has no target divisor, and a
self-adjoint similarity chosen after diagonalization would be noncanonical.
Thus
\[
\begin{aligned}
(A0,A1,A2,A3,A4)=(&\mathrm{A0\_WEAK\_ARITHMETIC\_RELATION},
\mathrm{A1\_FAIL},\\
&\mathrm{A2\_FAIL},\mathrm{A3\_FAIL},\mathrm{A4\_FORMAL\_HINT}).
\end{aligned}
\]
The scope literal is \texttt{NO\_BAD\_EULER\_OR\_ROOT\_NUMBER}.
The overall verdict is rejection and Route B is false.  No target tables,
arithmetic local data, Euler factors, root numbers, automorphy, target
functional equation, or Hilbert--P\'olya operator is claimed.  Nor is the
stochastic kernel promoted to an unweighted deterministic primitive-orbit
system on a changed phase space.

\begin{thebibliography}{2}
\bibitem{Holte97}
J.~M. Holte,
\newblock Carries, combinatorics, and an amazing matrix,
\newblock \emph{Amer. Math. Monthly} 104 (1997), 138--149,
\newblock \href{https://doi.org/10.1080/00029890.1997.11990612}{doi:10.1080/00029890.1997.11990612}.

\bibitem{DF09}
P.~Diaconis and J.~Fulman,
\newblock Carries, shuffling, and symmetric functions,
\newblock \emph{Adv. Appl. Math.} 43 (2009), 176--196,
\newblock \href{https://doi.org/10.1016/j.aam.2009.02.002}{doi:10.1016/j.aam.2009.02.002}.
\end{thebibliography}
\end{document}
